%!TEX root = ../main.tex
% APÉNDICE G — Protocolo de las sesiones moderadas con usuarios
% Fuente: docs/superpowers/specs/tesis/T11b-evaluacion-sin-piloto.md (§2)
% del repositorio del sistema. Resultados: capítulo 7 (sec:sesiones-moderadas).

\chapter{Protocolo de las sesiones moderadas}\label{apx:sesiones}

Este apéndice documenta el instrumento que reemplaza al piloto con
usuarios reales (sección~\ref{sec:sin-piloto}). Reproduce el protocolo
definido en la especificación de la evaluación del repositorio del
sistema (\path{docs/superpowers/specs/tesis/T11b-evaluacion-sin-piloto.md},
anexo~\ref{anx:repositorios}). Los resultados se informan en la
sección~\ref{sec:sesiones-moderadas}.

\section{Participantes y entorno}

\begin{itemize}
    \item Entre cinco y ocho personas \textbf{adultas}. Ningún menor de
          edad participa, ni como participante ni como beneficiario de
          una cuenta.
    \item Mezcla buscada, no obligatoria: al menos dos que hagan de
          tutor, al menos tres de alumno o adulto responsable y al menos
          una que viva en una provincia del interior (indicador I-03).
    \item Cuentas creadas durante la sesión con datos de prueba: el
          moderador entrega un \gls{dni} y una foto de \gls{dni} de
          prueba; no se suben documentos reales. Si el registro no acepta
          el documento de prueba, se usa una cuenta de los datos de prueba
          y la tarea de registro se anota como hecha por el moderador.
    \item Demostración pública con datos de prueba, pasarela de pagos en
          Modo Bypass (sin dinero real) y registro de evaluación encendido
          solo durante el período de las sesiones.
    \item Modalidad remota (videollamada con pantalla compartida) o
          presencial; unos 45 minutos por persona.
\end{itemize}

\section{Consentimiento informado}

Se lee y se firma antes de empezar. El texto es un borrador que el autor
revisa antes de usarlo; no contiene cláusulas legales que no se hayan
definido.
% TODO-AUTOR: revisar el texto del consentimiento, decidir si se graba pantalla o audio (y, en ese caso, pedir permiso aparte) y completar el contacto, el responsable de los datos y el plazo de conservación.

\begin{quote}
Te invitamos a probar \sistema{}, una plataforma de clases particulares
en línea que es parte de un trabajo final de carrera. Vas a usar una
versión de demostración: los tutores, las clases y los pagos son de
prueba, y no se cobra dinero real.

Durante unos 45 minutos te vamos a pedir que hagas algunas tareas
(registrarte, buscar un tutor, reservar, entrar a una clase de prueba,
calificar) y que digas en voz alta lo que pensás. Al final vas a
completar un cuestionario corto.

Qué registramos: si completaste cada tarea, cuánto tardaste, tus
respuestas al cuestionario y la provincia en la que vivís. La aplicación
además guarda, sin tu nombre a la vista, qué tutores te sugirió la
búsqueda y cuál abriste o reservaste, y tu respuesta a la pregunta de
recomendación (0 a 10). No guardamos el texto que escribas en la
búsqueda.

Para qué: los resultados agregados, sin nombres, se usan en el trabajo
final. Tu participación es voluntaria. Podés dejar de participar en
cualquier momento sin dar explicaciones y pedir que borremos tus datos
de la sesión.

$\square$ Leí lo anterior y acepto participar. Nombre: \rule{4cm}{0.4pt}
Fecha: \rule{2.5cm}{0.4pt}
\end{quote}

\section{Rol del moderador}

\begin{itemize}
    \item Lee cada tarea en voz alta tal como está escrita y no nombra
          botones.
    \item No ayuda, salvo que la persona esté trabada más de tres minutos
          en una tarea o pida ayuda dos veces; en ese caso anota la tarea
          como ``con ayuda'', que no cuenta como éxito.
    \item Antes de cada sesión verifica que el registro de evaluación esté
          encendido, que la pasarela esté en Modo Bypass y que los tutores
          de prueba tengan horarios libres en las horas siguientes.
\end{itemize}

\section{Guion de tareas}

\begin{table}[H]
    \centering
    \caption{Tareas del lado del alumno o adulto responsable}
    \label{tab:tareas-alumno}
    \small
    \begin{tabularx}{\textwidth}{l L L}
        \toprule
        \tb{\#} & \tb{Tarea, como se lee} & \tb{Éxito} \\
        \midrule
        T1 & ``Creá una cuenta en \sistema{} con los datos que te damos.'' & Llega al panel con la sesión iniciada \\
        T2 & ``Necesitás ayuda con [tema del sobre]. Buscá un tutor.'' & Ve una lista de tutores \\
        T3 & ``Elegí el tutor que te parezca mejor y mirá su perfil.'' & Abre el perfil de un tutor de la lista \\
        T4 & ``Reservá una clase de 30 minutos para hoy o mañana.'' & La reserva queda creada \\
        T5 & ``Pagá la clase.'' (Modo Bypass) & La reserva queda confirmada \\
        T6 & ``Entrá a la clase de prueba.'' (el moderador entra como tutor) & Ve el aula con su cámara \\
        T7 & ``Terminá la clase y calificá al tutor.'' & La calificación queda guardada \\
        T8 & Opcional, si presta un teléfono Android de gama media: ``Abrí esta dirección y tocá Empezar medición.'' & Copia el resultado de la página de diagnóstico del clasificador \\
        \bottomrule
    \end{tabularx}
\end{table}

\begin{table}[H]
    \centering
    \caption{Tareas del lado del tutor}
    \label{tab:tareas-tutor}
    \small
    \begin{tabularx}{\textwidth}{l L L}
        \toprule
        \tb{\#} & \tb{Tarea, como se lee} & \tb{Éxito} \\
        \midrule
        U1 & ``Creá tu cuenta de tutor.'' & Llega al panel de tutor \\
        U2 & ``Cargá las materias que podés enseñar.'' & Guarda al menos un tema \\
        U3 & ``Publicá tus horarios para esta semana.'' & Guarda al menos una franja \\
        U4 & ``Fijá tu tarifa por hora.'' & Tarifa guardada, respetando el piso \\
        U5 & ``Entrá a la clase que te reservaron y terminala.'' & La sesión queda finalizada \\
        U6 & ``Calificá al alumno.'' & Calificación guardada \\
        \bottomrule
    \end{tabularx}
\end{table}

Después de la primera clase finalizada o calificada, la aplicación
ofrece la pregunta de recomendación. El moderador no la menciona; si la
persona la cierra sin responder, se le pregunta en el cuestionario.

\section{Métricas y metas}

\begin{table}[H]
    \centering
    \caption{Métricas de las sesiones moderadas}
    \label{tab:metricas-sesiones}
    \small
    \begin{tabularx}{\textwidth}{L L l}
        \toprule
        \tb{Métrica} & \tb{Cómo se obtiene} & \tb{Meta} \\
        \midrule
        Éxito por tarea & Éxitos sin ayuda sobre participantes que la intentaron & $\geq 80\,\%$ en T2--T7 y U2--U5 \\
        Tiempo por tarea & Cronómetro del moderador, de la lectura al criterio de éxito & Mediana y máximo, sin meta \\
        SUS & Cuestionario posterior, puntaje de 0 a 100 & Promedio $\geq 68$ \\
        \gls{nps} de intención & En la aplicación o en el cuestionario & Positivo \\
        I-01b & Panel de evaluación del administrador & $> 60\,\%$ \\
        I-02 & Calificaciones del alumno al tutor de esas clases & $\geq 4{,}0$ \\
        I-03 & Provincia declarada en el cuestionario & $\geq 1$ clase entre provincias, una del interior \\
        \bottomrule
    \end{tabularx}
\end{table}

I-01b se calcula como la proporción de búsquedas con al menos un
resultado en las que la persona abrió el perfil o reservó a un tutor de
los cinco primeros. El \gls{nps} es el porcentaje de promotores (9 y 10)
menos el de detractores (0 a 6). La meta de SUS es el promedio de
referencia de la escala \citep{sauro2011}.

\section{Cuestionario posterior}

\textbf{Datos:} provincia en la que vive (una de las 24 jurisdicciones),
rol que probó y si alguna vez tomó o dio clases particulares en línea.

\textbf{SUS} \citep{brooke1996}, en castellano, con una escala de 1 (muy
en desacuerdo) a 5 (muy de acuerdo):

\begin{enumerate}
    \item Creo que me gustaría usar \sistema{} con frecuencia.
    \item Encontré \sistema{} innecesariamente complejo.
    \item Me pareció que \sistema{} era fácil de usar.
    \item Creo que necesitaría ayuda de una persona con conocimientos
          técnicos para poder usar \sistema{}.
    \item Encontré que las distintas funciones de \sistema{} estaban bien
          integradas.
    \item Me pareció que había demasiadas inconsistencias en \sistema{}.
    \item Imagino que la mayoría de las personas aprenderían a usar
          \sistema{} muy rápido.
    \item Encontré \sistema{} muy incómodo de usar.
    \item Me sentí muy seguro o segura usando \sistema{}.
    \item Necesité aprender muchas cosas antes de poder empezar a usar
          \sistema{}.
\end{enumerate}

Cálculo: en los ítems impares se resta 1 a la respuesta; en los pares,
la respuesta se resta de 5; la suma se multiplica por 2,5.

\textbf{Recomendación} (solo si no la respondió en la aplicación): ``¿Qué
tan probable es que recomiendes \sistema{} a un amigo o familiar? (0 =
nada probable, 10 = muy probable)''.

\textbf{Preguntas abiertas:} ``¿Qué fue lo más difícil?'' y ``¿Qué
cambiarías primero?''.

\section{Planilla de resultados}

Una fila por participante, identificada como P01 a P08, sin nombres, con
el rol, la provincia, si es del interior, el resultado (sí, con ayuda,
no o no aplica) y el tiempo en segundos de cada tarea, las diez
respuestas del SUS y su total, el \gls{nps} y su fuente, el participante
con el que tuvo la clase (para I-03) y observaciones. La hoja de resumen
calcula la tasa de éxito y el tiempo mediano y máximo por tarea, el SUS
promedio y su desvío, el \gls{nps} e I-01b, I-02 e I-03 con su fecha.
